This subsection is generally organized in order of increasing strength of hypotheses and some results are proved in greater generality than stated in \Cref{thm:synthetic-Adams}. Before we begin we will need to recall more material from \cite{Pstragowski}.

\begin{rec}
  In \cite[Section 4.2, Lemma 4.29 and Proposition 4.35]{Pstragowski},
  Pstragowski introduces a $t$-structure on $\Syn_E$ which satisfies the following properties:
  \begin{enumerate}
  \item The heart, $\mathrm{Syn}_E ^{\heartsuit}$, is equivalant to the abelian category of $E_* E$-comodules.
  \item Given a spectrum $X$, the $0^{\mathrm{th}}$ homotopy object of $\nu X$ with respect to this $t$-structure is naturally equivalent to $E_*X$ as an $E_*E$-comodule and naturally equivalent to $C\tau \otimes \nu(X)$ as an object of $\Syn_E$.
  \item If we let $Y(-)$ denote the right adjoint to inverting $\tau$,
    then we have a natural equivalence between the connective cover of $Y(X)$ and $\nu X$.  
  \item This $t$-strucutre is right complete and compatible with filtered colimits.
  \end{enumerate}
\end{rec}

\begin{ntn}
  In view of (2) we will refer to the $t$-structure introduced above as the \emph{homological $t$-structure} on $\Syn_E$.
  For notational brevity we will use subscripts to denote truncation with respect to this $t$-structure, so that $A_{\geq n}$ refers to the $n$-connective cover of a synthetic spectrum $A$ in the homological $t$-structure.  
\end{ntn}

\begin{cnv}
For the remainder of this subsection $X$ will denote a spectrum.
\end{cnv}
Our analysis of the relation between the bigraded homotopy groups of $\nu X$ and the $E$-based Adams spectral sequence for $X$ will hinge on an understanding of the $\nu E$-based Adams spectral sequence for $\nu X$. We begin by considering the canonical $E$-Adams tower for $X$, constructed below:

\begin{center}
  \begin{tikzcd}
    \cdots \ar[r] &
    X_2 \ar[r, "f_2"] \ar[d, "i_2"] &
    X_1 \ar[r, "f_1"] \ar[d, "i_1"] &
    X_0 \ar[r, equal] \ar[d, "i_0"] &
    X. \\
    & E \otimes X_2 &
    E \otimes X_1 &
    E \otimes X_0 &
  \end{tikzcd}
\end{center}
Each of the $f_i$ is zero on $E$-homology. Therefore, using \Cref{rmk:strong-monoidal} and \Cref{lemm:syn-cof}, we may identify the canonical $\nu E$-Adams tower of $\nu X$ as
\begin{center}
  \begin{tikzcd}
    &
    \Sigma^{0,2} \nu X_3 \ar[d, "\nu(f_3)"] \ar[dl, "\tau"'] &
    \Sigma^{0,1} \nu X_2 \ar[d, "\nu(f_2)"] \ar[dl, "\tau"'] &
    \nu X_1 \ar[d, "\nu(f_1)"] \ar[dl, "\tau"'] & \\
    \cdots \ar[r] &
    \Sigma^{0,2} \nu X_2 \ar[r, "\wt{f_2}"] \ar[d, "\nu(i_2)"] &
    \Sigma^{0,1} \nu X_1 \ar[r, "\wt{f_1}"] \ar[d, "\nu(i_1)"] &
    \nu X_0 \ar[r, equal] \ar[d, "\nu(i_0)"] &
    \nu X. \\
    &
    \nu E \otimes \Sigma^{0,2} \nu X_2 &
    \nu E \otimes \Sigma^{0,1} \nu X_1 &
    \nu E \otimes \nu X_0 &
  \end{tikzcd}
\end{center}

%% s is the filtration
%% k is the stem
%% w is the weight

\begin{ntn}
The above tower gives rise to a spectral sequence
 $$ \mathrm{E}^{s,k,w}_1 \coloneqq \pi_{k, k+w} \left( \nu E \otimes \Sigma^{0,s} \nu X_s \right) \Longrightarrow \pi_{k, k+w}(\nu X), $$
 with differentials of tridegree $(r, -1, 1)$. Note that multiplication by $\tau$ lowers the $w$ grading by $1$ but preserves the $s$ and $k$ gradings.
  We use the notation $\mathrm{E}^{s,k,w}_r$ for page $r$ of this spectral sequence.

  Analogously, we use $\mathrm{E}^{s,k}_r$ to refer to the groups in the $E$-Adams spectral sequence for $X$
 $$ \mathrm{E}^{s,k}_1 \coloneqq \pi_k \left( E \otimes X_s \right) \Longrightarrow \pi_k(X), $$
with differentials of bidegree $(r, -1)$. \footnote{Our grading choices do not agree with the usual conventions for Adams spectral sequences. However, we prefer them because each of the indices has a clear interpretation: $k$ is the topological degree, $w$ is the weight and $s$ is the filtration.}
\end{ntn}
 Note that inverting $\tau$ determines a map of spectral sequences
 $$ \mathrm{E}_r^{s,k,w} \to \mathrm{E}_r^{s,k}. $$

 \begin{ntn}
 Let $\mathrm{B}_r^{s,k}$ denote the subgroup of $\mathrm{E}_2^{s,k}$ generated by the images of the differentials $d_2$ through $d_r$. Let $\mathrm{Z}_r^{s,k}$ denote the larger subgroup of $\mathrm{E}_2^{s,k}$ given by those classes on which $d_2$ through $d_r$ vanish. Then, $\mathrm{E}_{r+1}^{s,k} \cong \mathrm{Z}_r^{s,k}/\mathrm{B}_r^{s,k}$.
 \end{ntn}

\begin{thm} \label{thm:synth-ASS}
    The $\nu E$-based Adams spectral sequence for $\nu X$ is determined by the $E$-based Adams spectral sequence for $X$ in the following way:
    \begin{enumerate}
    \item $ \mathrm{E}^{s,k,w}_1 \cong \mathrm{E}^{s,k}_1 \otimes \ZZ[\tau]$,
      where $\mathrm{E}^{s,k}_1$ is considered to be in tridegree $(s,k,s)$.
    \item $ \mathrm{E}^{s,k,w}_2 \cong \mathrm{E}^{s,k}_2 \otimes \ZZ[\tau]$,
      where $\mathrm{E}^{s,k}_2$ is considered to be in tridegree $(s,k,s)$.
    \item Given a differential $d_{r, \mathrm{top}} (x) = y$, there is a differential $d_r (x) = \tau^{r-1} y$. Moreover, all differentials arise in this way.
    \end{enumerate}
\end{thm}

\begin{proof}
  The proof is very similar to that of \cite[Lemma 15]{HKO}.
  Statement (1) follows from \cite[Proposition 4.21]{Pstragowski}.
  Statement (2) follows from statement (3).
  We now prove statement (3) by induction. Suppose that we have proved the statement through the $\mathrm{E}_r$-page.  To prove it for the $\mathrm{E}_{r+1}$-page, we calculate the differential
  $$ d_r : \mathrm{E}^{s,k,w}_r \to \mathrm{E}_r^{s+r,k-1,w+1}. $$
  Note that, by the inductive hypothesis, the $\mathrm{E}_r$-page in every tridegree which can be the target of a $d_r$ differential consists of $\tau$-torsion free elements.  On the other hand, upon inverting $\tau$ we must obtain the differential
$$\tau^{-1} d_r = d_{r,\mathrm{top}}: \mathrm{E}^{s,k}_{r} \to \mathrm{E}^{s+r,k-1}_{r},$$
which determines the $d_r$ differential by the above.
\end{proof}

As a corollary of this description of the $\nu E$-Adams spectral sequence, we obtain the following more explicit statement.

\begin{cor} \label{cor:synth-ASS}
  For $2 \leq r \leq \infty$, there are natural isomorphisms:
  \begin{enumerate}
  \item $\mathrm{E}_r^{s,k,w} \cong 0$ for $w > s$.
  \item $\mathrm{E}_r^{s,k,w} \cong \mathrm{Z}_{r-1}^{s,k}$ for $w = s$.
  \item $\mathrm{E}_r^{s,k,w} \cong \mathrm{Z}_{r-1}^{s,k}/\mathrm{B}_{s-w+1}^{s,k}$ for $s-r+1 < w < s$.
  \item $\mathrm{E}_r^{s,k,w} \cong \mathrm{E}_r^{s,k}$ for $w \leq s-r+1$ and $w \leq 0$.
  \end{enumerate}
    In particular, the map $\mathrm{E}_r ^{s,k,w} \to \mathrm{E}_r ^{s,k,w-1}$ induced by multiplication by $\tau$ is surjective for $w \leq s$.
\end{cor}

Our next order of business will be to determine the $\nu E$-based Adams spectral sequence for $C \tau^p \otimes \nu X$.

\begin{ntn}
    We use the notation ${}^p \mathrm{E}_r ^{s,k,w}$ to denote the groups on page $r$ of the $\nu E$-based Adams spectral sequence
    \[ {}^p\mathrm{E}_1^{s,k,w} \coloneqq \pi_{k,k+w}( \nu E \otimes C\tau^p \otimes \Sigma^{0,s}\nu X_s ) \Longrightarrow \pi_{k,k+w} (\nu X),\]
  and similarly for the later pages.
\end{ntn}

%By considering the map $\nu X \to C\tau^p \otimes \nu X$ we can deduce the structure of the target's $\nu E$-based Adams spectral sequence as a corollary to \Cref{cor:synth-ASS}. 
%
\begin{cor} \label{cor:synth-ctau-ASS}
  For $p \geq 1$ and $2 \leq r \leq \infty$, there are natural isomorphisms:
  \begin{enumerate}
  \item ${}^p\mathrm{E}_r^{s,k,w} \cong 0$ for $w > s$.
  \item ${}^p\mathrm{E}_\infty^{s,k,w} \cong \mathrm{Z}_{p-s+w}^{s,k}/\mathrm{B}_{s-w+1}^{s,k}$ for $s \geq w > s-p$.
  \item ${}^p\mathrm{E}_r^{s,k,w} \cong 0$ for $s-p \geq w$.
  \end{enumerate}
\end{cor}

\begin{proof}
    This follows from considering the map of $\nu E$-based Adams spectral sequences induced by the map $\nu X \to C\tau^p \otimes \nu X$. \todo{Maybe elaborate on this.}
\end{proof}

In order to use the theorem and corollaries we have just proved we will need to make a digression and discuss completeness and convergence.

\begin{dfn}
  We say that a synthetic spectrum $A$ is \emph{$\tau$-complete} if the $\tau$-Bockstein tower of $A$ is convergent: that is, if the canonical map
  \[A \to \varprojlim_n C\tau^n \otimes A\] is an equivalence.
\end{dfn}

\begin{prop} \label{lemm:Elocal-taucomplete} \label{lemm:easy-e-compton}
  The following are equivalent:
  \begin{enumerate}
  \item $X$ is $E$-nilpotent complete.
  \item $\nu X$ is $\nu E$-nilpotent complete.
  \item $\nu X$ is $\tau$-complete.
  \end{enumerate}
\end{prop}

The proof of \Cref{lemm:Elocal-taucomplete} will rely on the following two lemmas.

\begin{lem} \label{lemm:comp1}
  The synthetic spectrum $C\tau^p \otimes \nu X$ is $\nu E$-nilpotent complete.
\end{lem}

\begin{proof}
    By induction on $p$ via the the Bockstein sequences
    \[\Sigma^{0,-1} C\tau^{p-1} \otimes \nu X \to C \tau^p \otimes \nu X \to C \tau \otimes \nu X,\]
  we see that it suffices to prove the lemma for $p=1$.
  Tensoring the canonical Adams resolution for $\nu X$ with $C\tau$, we obtain an Adams resolution of $C\tau \otimes \nu X$.
  Using \cite[Lemma 4.29]{Pstragowski} repeatedly, we learn that $C\tau \otimes \Sigma^{0,s}\nu X$ is $(-s)$-coconnective in the homological $t$-structure. Thus, the inverse limit of this Adams resolution for $\nu X$ is trivial because this $t$-structure is right complete.
\end{proof}

\begin{lem} \label{lemm:comp2}
  The synthetic spectrum $\nu E \otimes \nu X$ is $\tau$-complete.
\end{lem}

\begin{proof}
  In order to show that $\nu E \otimes \nu X$ is $\tau$-complete we will show that the inverse limit under iterated multiplication by $\tau$ is trivial.
  From \Cref{rmk:dualizable-generators} it suffices to check triviality on maps in from suspensions of finite projectives.
  Pick a finite $E_*$-projective $P$ and an integer $k$. Using \Cref{rmk:strong-monoidal} and the dualizability statement from \Cref{rmk:dualizable-generators}, we obtain an equivalence
  $$ \Hom(\Sigma^k \nu P, \varprojlim_\tau \Sigma^{0,-s}\nu E \otimes \nu X) \simeq \varprojlim_\tau \Hom(\Ss^{k,s}, \nu (E \otimes DP \otimes X)). $$
    %% &\simeq \hom(\Sigma^k \nu P, \varprojlim_\tau \Sigma^{0,-s}\nu (E \otimes X)) \\
    %% &\simeq\varprojlim_\tau \hom(\Sigma^k \nu P, \Sigma^{0,-s}\nu (E \otimes X)) 
    %% &\simeq\varprojlim_\tau \hom(\Ss^{k,s}, \nu DP \otimes \nu (E \otimes X)) \\
    %% &\simeq 
  Note that the spaces in the inverse limit on the right hand side are each $(s-k)$-connective by \cite[Proposition 4.21]{Pstragowski}. Therefore as $s \to \infty$ the right hand side becomes infinitely connective and thereby trivial.
\end{proof}

\begin{proof}[Proof of \Cref{lemm:easy-e-compton}]
  First we show that $E$-nilpotent completeness is equivalent to $\nu E$-nilpotent completeness.
  Using \cite[Lemma 4.29 and Proposition 4.35]{Pstragowski}, we may rewrite the canonical $\nu E$-Adams tower for $\nu X$ as
  $$ \cdots \to Y(X_2)_{\geq -2} \xrightarrow{\wt{f_2}} Y(X_1)_{\geq -1} \xrightarrow{\wt{f_1}} Y(X_0)_{\geq 0},$$
	where
	$\cdots \to X_2 \to X_1 \to X_0$
	is the canonical Adams tower for $X$.
  Inverting $\tau$ on the $\nu E$-Adams tower recovers the image of the $E$-Adams tower under $Y(-)$, and there is a fiber sequence
  $$ \left( \varprojlim_s Y(X_s)_{\geq -s} \right) \to \left( \varprojlim_s Y(X_s) \right) \to \left( \varprojlim_s Y(X_s)_{\leq -s} \right). $$
  The right hand term vanishes because the homological $t$-structure is right complete \cite[Proposition 4.16]{Pstragowski}. Furthermore, since $Y$ is a right adjoint, we may pull the inverse limit inside the functor $Y$ in the middle term. Thus, we obtain an equivalence
  $$ \left( \varprojlim_s Y(X_s)_{\geq -s} \right) \simeq Y \left( \varprojlim_s X_s \right).$$
  From the fact that $Y$ is fully faithful we conclude that the left hand side vanishes if and only if the inverse limit of the $E$-Adams tower for $X$ vanishes.
  
  Next we show that $\nu E$-nilpotent completeness is equivalent to $\tau$-completeness. Consider the following diagram:
%\begin{center}
$$  \begin{tikzcd}[column sep=small]
    \varprojlim_s \varprojlim_\tau \Sigma^{0,s-p} \nu X_s \ar[r] & \cdots \ar[r] &
    \varprojlim_\tau \Sigma^{0,1-p} \nu X_1 \ar[r, "\wt{f_1}"] \ar[d, "\nu(i_1)"] &
    \varprojlim_\tau \Sigma^{0,-p}\nu X \ar[d, "\nu(i_0)"] \\
    & &
    \varprojlim_\tau \Sigma^{0,1-p} (\nu E \otimes \nu X_1) &
    \varprojlim_\tau \Sigma^{0,-p} (\nu E \otimes \nu X). &
			\end{tikzcd} $$
%\end{center}
Here, the limits over $\tau$ refer to limits, as $p$ varies, under multiplication by $\tau$ maps.
Using \Cref{lemm:comp2}, each object on the second row vanishes. We obtain an equivalence
$$ \varprojlim_s \varprojlim_\tau \Sigma^{0,s-p}\nu X_s \simeq \varprojlim_\tau \Sigma^{0,-p}X. $$
Dually, using \Cref{lemm:comp1}, we learn that
$$ \varprojlim_\tau \varprojlim_s \Sigma^{0,s-p}\nu X_s \simeq \varprojlim_s \Sigma^{0,s}\nu X_s. $$
Together these equalities finish the proof.
\end{proof}

%% \begin{dfn}
%%   We will say that a $\nu E$-nilpotent complete synthetic spectrum $Y$ has \emph{strongly convergent} $\nu E$-based Adams spectral sequence if the $\nu E$-Adams filtration on its bigraded homotopy groups is complete and Hausdorff with associated graded isomorphic to the appropriate $\mathrm{E}_\infty$-page.

%%   We also make a similar definiton for the $\tau$-Bockstein spectral sequence.
%% \end{dfn}

% In particular, we conclude that the $\nu E$-based Adams spectral sequence for $\nu X$ is conditionally convergent whenever $X$ is $E$-nilpotent complete. Using the structure of the $\nu E$-based Adams spectral sequence for $\nu X$ from \Cref{thm:synth-ASS}, we will show in \Cref{cor:strong-conv} that that $\nu E$-based Adams spectral sequence for $\nu X$ converges strongly whenever the $E$-based Adams spectral sequence for $X$ does.
% In order to directly relate the final page of an Adams spectral sequence to homotopy groups we need slightly more than just completeness.
% The analog of \Cref{lemm:Elocal-taucomplete} then holds for these stronger notions of completeness as well and as above we will defer the proof to ???.

We are now ready to prove the first part of \Cref{thm:synthetic-Adams}.

\begin{proof}[Proof of \Cref{thm:synthetic-Adams}(1)]
  Since $X$ is $E$-nilpotent complete, it follows from \Cref{lemm:easy-e-compton} that $C\tau^r \otimes \nu X$ is $\nu E$-nilpotent complete and $\tau$-complete. From \Cref{cor:synth-ctau-ASS} we can read off that the $\nu E$-based Adams spectral sequence for $C\tau^r \otimes \nu X$ converges strongly. Further, we can directly read off that (1a) and (1b) are equivalent. Clearly (1c) implies (1b). We will now prove (1c) assuming (1b).  If $d_{r+1}(x) = 0$, then we may finish by (1a), so we assume otherwise.

  We will prove (1c) by working directly with the cofiber sequence of Adams towers associated to the relevant Bockstein sequence. Before we begin, we fix some notation,
  $$ {}^r\mathrm{D}^{s,k,w} \coloneqq \pi_{k,k+w}( C\tau^r \otimes \Sigma^{0,s}\nu X_s ). $$
	Now, consider the following diagram of exact sequences, obtained by applying homotopy groups to the smash product of the two extended cofiber sequences
\[\Sigma^{-1,0} C\tau \to C\tau^{r+1} \to C\tau^{r} \to C\tau, \text{ and}\]
\[\Sigma^{0,s+1} \nu X_{s+1} \to \Sigma^{0,s} \nu X_{s} \to \nu E \otimes \Sigma^{0,s} \nu X_{s} \to \Sigma^{1,s+1} \nu X_{s+1}.\]
  \begin{center}
    \begin{tikzcd}
      {}^1\mathrm{D}^{s+1,k,s+r} \ar[r] \ar[d, "f"] &
      {}^{r+1}\mathrm{D}^{s+1,k,s} \ar[r] \ar[d, "f"] &
      {}^r\mathrm{D}^{s+1,k,s} \ar[r] \ar[d, "f"] &
      {}^1\mathrm{D}^{s+1,k-1,s+r+1} \ar[d, "f"] \\
      {}^1\mathrm{D}^{s,k,s+r} \ar[r] \ar[d, "i"] &
      {}^{r+1}\mathrm{D}^{s,k,s} \ar[r] \ar[d, "i"] &
      {}^r\mathrm{D}^{s,k,s} \ar[r] \ar[d, "i"] &
      {}^1\mathrm{D}^{s,k-1,s+r+1} \ar[d, "i"] \\
      0 \ar[r] \ar[d, "\partial"] &
      {}^{r+1}\mathrm{E}_1^{s,k,s} \ar[r] \ar[d, "\partial_{r+1}"] &
      {}^r\mathrm{E}_1^{s,k,s} \ar[r] \ar[d, "\partial_r"] &
      0 \ar[d, "\partial"] \\
      {}^1\mathrm{D}^{s+1,k-1,s+r+1} \ar[r] &
      {}^{r+1}\mathrm{D}^{s+1,k-1,s+1} \ar[r] &
      {}^r\mathrm{D}^{s+1,k-1,s+1} \ar[r] &
      {}^1\mathrm{D}^{s+1,k-2,s+r+2} 
    \end{tikzcd}
  \end{center}
	
  In this diagram we may pick a representative of $x$ in ${}^r\mathrm{E}_1^{s,k,s} = {}^{r+1}\mathrm{E}_1^{s,k,s}$ (which we will also denote $x$).
  Let $y = d_r(x)$ denote the target of the relevant differential in the $E$-based Adams spectral sequence and consider $y$ as an element of ${}^1\mathrm{D}^{s+1,k-1,s+r+1}$.
  We claim that there exists a $y' \in {}^{r+1}\mathrm{D}^{s+1,k-1,s+r+1}$ such that $\partial_{r+1} (x) = \tau^r y'$ and $y$ maps to $\tau^r y'$.
  Indeed, this follows from \Cref{thm:synth-ASS} and the fact that $\tau^r$ as an endomorphism of $C\tau^{r+1}$ factors through $C\tau$.

  Then, by standard manipulations of exact sequences arising from smash products of cofiber sequences (as in, e.g., \cite[Lemma 9.3.2]{AndrewsMiller}), there is some $\wt{x}$ in ${}^r\mathrm{D}^{s,k,s}$ which maps to both $x$ in ${}^r\mathrm{E}_1^{s,k,s}$ and $-f(y)$ in ${}^1\mathrm{D}^{s,k-1,s+r+1}$. The image of $\wt{x}$ along the map
  $$ {}^r\mathrm{D}^{s,k,s} \to \pi_{k,k+s}(C\tau^r \otimes \nu X) $$
  is the desired class.
\end{proof}

\begin{prop} \label{lemm:strong-etau-complete} \label{cor:strong-conv}
  Let $X$ denote an $E$-nilpotent complete spectrum.
  Then the following are equivalent:
  \begin{enumerate}
  \item The $E$-based Adams spectral sequence for $X$ converges strongly.
  \item The $\nu E$-based Adams spectral sequence for $\nu X$ converges strongly.
  \item The $\tau$-Bockstein spectral sequence for $\nu X$ converges strongly.
  \end{enumerate}
\end{prop}

In order to prove \Cref{lemm:strong-etau-complete} we recall the following theorem of Boardman, which provides a useful characterization of strong convergence.

\begin{thm}[{\cite[Theorem 7.3]{Boardman}}] \label{thm:strong-conv}
    Given an $E$-nilpotent complete spectrum $X$, the following two conditions are equivalent:
    \begin{itemize}
        \item The $E$-based Adams spectral sequence of $X$ converges strongly.
        \item $\varprojlim_{r}^{1} \mathrm{E}^{s,t}_r (X) = 0$ for pair of integers $s$ and $t$.
    \end{itemize}
		Analogous $\mathrm{lim}^1$ conditions determine strong convergence of $\nu E$-based Adams spectral sequences and $\tau$-Bockstein spectral sequences.
\end{thm}

Note that the second bullet point in the above theorem makes sense because $\mathrm{E}^{s,t} _{r+1} \subseteq \mathrm{E}^{s,t} _r$ as soon as $r>s$.

\begin{proof}[Proof of \Cref{lemm:strong-etau-complete}]
  Since $X$ is $E$-nilpotent complete, it follows from \Cref{lemm:easy-e-compton} that $\nu X$ is $\nu E$-nilpotent complete and $\tau$-complete.

  Using \Cref{thm:strong-conv}, to prove that (1) is equivalent to (2) it suffices to show that $\varprojlim_r \mathrm{E}^{s,k}_{r} = 0$ if and only if $\varprojlim_r \mathrm{E}^{s,k,w}_{r} = 0$. 
  In fact, these groups are isomorphic:
    \begin{align*}
      {\varprojlim_{r}}^1 \mathrm{E}^{s,k} _{r}
      \cong {\varprojlim_{r}}^1 \mathrm{Z}^{s,k}_{r}/\mathrm{B}_{s}^{s,k} 
      \cong {\varprojlim_{r}}^1 \mathrm{Z}^{s,k}_{r} 
      \cong {\varprojlim_{r}}^1 \mathrm{Z}^{s,k}_{r}/\mathrm{B}_{s-w+1}^{s,k} 
      \cong {\varprojlim_{r}}^1 \mathrm{E}^{s,k,w}_{r}.
    \end{align*}

We next prove the equivalence of the second and third conditions.  Let $\beta_r^{s,k,w}$ denote the groups in the $\tau$-Bockstein spectral sequence indexed so that the spectral sequence takes the form
    $$\beta_1^{s,k,w} \cong \pi_{k,k+w}(\Sigma^{0,-s} C\tau \otimes \nu X) \Longrightarrow \pi_{k,k+w} (\nu X).$$
Combining \Cref{thm:synthetic-Adams}(1) and \Cref{cor:synth-ASS}, we learn that
	$$\beta_r^{s,k,w} \cong \mathrm{E}_{r+1}^{w+s,k,w}.$$
Boardman's theorem applies since both of these spectral sequences are conditionally convergent.  Since the spectral sequences are furthermore isomorphic, up to reindexing, one converges strongly if and only if the other does.
%\begin{comment}
%    Next, we prove that the first and third conditions are equivalent.
%    Let $\beta_r^{s,k,w}$ denote the $\tau$-Bockstein spectral sequence indexed so that
%    $$\beta_1^{s,k,w} \cong \pi_{k,k+w}(C\tau \otimes \nu X).$$
%    Using the facts that
%    \begin{itemize}
%    \item the spectral sequence is concentrated in the region with $s \geq 0$,
%    \item $d_r$ differentials increase $s$ by $r$, and
%    \item multiplication by $\tau$ is surjective on each page,
%    \end{itemize}
%    we learn that the $\tau$-Bockstein spectral sequence converges strongly if and only if the $\lim^1$ terms of the groups on the $s=0$ plane vanish. To finish, it suffices to observe by \Cref{thm:synthetic-Adams}(1a,b) that there are isomorphisms
%    \[ \beta_r^{0,k,w} \cong Z_r^{w,k} \cong E_{r+1}^{w,k,w}. \qedhere \]
%\end{comment}
\end{proof}

\begin{ntn}
    We let $\mathrm{F}^s \pi_{k,k+w} (\nu X) \subseteq \pi_{k,k+w} (\nu X)$ denote the $\nu E$-Adams filtration, and we let $\mathrm{F}^{s}_{\tau} \pi_{k,k+w}(\nu X)$ denote the $\tau$-Bockstein filtration.
\end{ntn}

\begin{cor} \label{cor:tau-surj}
  Suppose $X$ is $E$-nilpotent complete and that its $E$-based Adams spectral sequence converges strongly. Then
    %$$\mathrm{F}^{s} \pi_{k,k+w}(\nu X) = \mathrm{F}^{0}_{\tau} \pi_{k,k+w}(\nu X) = \pi_{k,k+w}(\nu X).$$
	%If $s>w$, then
    $$\mathrm{F}^{s} \pi_{k,k+w}(\nu X) = \mathrm{F}^{s-w}_{\tau} \pi_{k,k+w}(\nu X),$$
    where for $k < 0$ we set
    \[\mathrm{F}^{k} _{\tau} \pi_{k,k+w} (\nu X) = \pi_{k,k+w} (\nu X).\]
    In particular, for $w \leq s$, the map
  $$ \mathrm{F}^s\pi_{k,k+w}(\nu X) \xrightarrow{\cdot \tau} \mathrm{F}^s\pi_{k,k+w-1}(\nu X) $$
  is surjective.
	
	\begin{comment}
	for $w \leq s$, the map
  $$ \mathrm{F}^s\pi_{k,k+w}(\nu X) \xrightarrow{\cdot \tau} \mathrm{F}^s\pi_{k,k+w-1}(\nu X) $$
  is surjective.  Furthermore, for $w \leq s$, $\mathrm{F}^{s + \bullet}\pi_{k,k+w}(\nu X)$ coincides with the $\tau$-adic filtration on $\mathrm{F}^{s + \bullet}\pi_{k,k+w}(\nu X)$.
	\end{comment}
\end{cor}

\begin{proof}
  By \Cref{cor:strong-conv}, the $\nu E$-based Adams spectral sequence for $\nu X$ converges strongly.
  %In particular, this means that
    %\begin{enumerate}
    %\item $\mathrm{F}^s \pi_{k,k+w}( \nu X )/ \mathrm{F}^{s+1} \pi_{k,k+w}( \nu X ) \cong \mathrm{E}^{s,k,w} _\infty$.
    %\item The filtration $\mathrm{F}^{s} \pi_{k,k+w}( \nu X )$ is complete and Hausdorff.
    %\end{enumerate}

    Now, the inclusion
    $$\mathrm{F}^{s-w} _{\tau} \pi_{k,k+w}(\nu X) \subseteq \mathrm{F}^{s} \pi_{k,k+w}(\nu X)$$
    follows from a downward induction on $w$, starting from \Cref{cor:synth-ASS}(1), which implies the desired result for $w \geq s$.
    On the other hand, to see that 
    $$\mathrm{F}^{s} \pi_{k,k+w}(\nu X) \subseteq \mathrm{F}^{s-w}_{\tau} \pi_{k,k+w}(\nu X)$$
    for all $s$, it suffices by strong convergence to show that, whenever $w \leq s$, multiplication by $\tau$ is surjective as a map $\mathrm{E}_{\infty} ^{s,k,w-1} \to \mathrm{E}_{\infty} ^{s,k,w}$. This is a consequence of \Cref{cor:synth-ASS}.
%
%    whenever $w \leq s$ multiplication by $\tau$ is surjective on associated graded. The identification of the filtration follows from the surjectivity statement and \Cref{cor:synth-ASS}(1).
\end{proof}

\begin{proof}[Proof of \Cref{thm:synthetic-Adams}(2)-(4)]
    We begin by noting that \Cref{cor:strong-conv} implies that the $\nu E$-based Adams spectral sequence for $\nu X$ converges strongly.
    Recall that this means that:

    \begin{enumerate}
    \item $\mathrm{F}^s \pi_{k,k+w}( \nu X )/ \mathrm{F}^{s+1} \pi_{k,k+w}( \nu X ) \cong \mathrm{E}^{s,k,w} _\infty$.
    \item The filtration $\mathrm{F}^{\bullet} \pi_{k,k+w}( \nu X )$ is complete and Hausdorff.
    \end{enumerate}

    Now, we examine the reduction map
    $$ \nu X \to C\tau \otimes \nu X $$
    through the $\nu E$-based Adams spectral sequence.
    As discussed in \Cref{cor:synth-ctau-ASS}, the $\nu E$-based Adams spectral sequence for $C\tau \otimes \nu X$ has $\mathrm{E}_2$-term given by
    $$ {}^1\mathrm{E}_2^{s,k,w} \cong \begin{cases} \mathrm{E}^{s,k}_{2}, & \text{if } s=w \\ 0, & \text{otherwise.} \end{cases} $$
    It follows that the spectral sequence collapses at the $\mathrm{E}_2$-page, that there is no space for extension problems and that the map
    $$ \mathrm{Z}_\infty^{s,k} \cong \mathrm{E}_\infty^{s,k,s} \to {}^1\mathrm{E}_\infty^{s,k,s} \cong \mathrm{E}_2^{s,k} $$
    is just the usual inclusion. This produces a factorization
    $$\pi_{k,k+s}(\nu X)=F^{s} \pi_{k,k+s}(\nu X) \twoheadrightarrow \mathrm{E}_\infty^{s,k,s} \cong \mathrm{Z}_\infty^{s,k} \subseteq \mathrm{E}_2^{s,k} \cong \pi_{k,k+s}(C\tau \otimes \nu X). $$
    The surjectivity of the first map implies that we can always pick an $\wt{x}$.

    (2a) On the associated graded, multiplication by $\tau^{r-1}$ can be identified with
    $$ \mathrm{E}_\infty^{s,k,s} \cong \mathrm{Z}_\infty^{s,k}(X) \to \mathrm{Z}_\infty^{s,k}(X) / \mathrm{B}_{r}^{s,k}(X). $$
    Therefore, as long as $x$ survives to the $\mathrm{E}_{r+1}$-page, any lift $\wt{x}$ will have $\tau^{r-1}\wt{x} \neq 0$.
    
    (2b) It suffices to note that the $\nu E$-based Adams spectral sequence for $\nu X$ is sent to the $E$-based Adams spectral sequence for $X$ under $\tau^{-1}$ and that the induced map
    $$ \mathrm{E}_\infty^{s,k,s} \cong \mathrm{Z}_\infty^{s,k}(X) \to \mathrm{E}_\infty^{s,k} $$
    is just the usual projection.

    (3a) For this we now suppose that $x \in \mathrm{B}_{r+1}^{s,k}$ and consider the following diagram:
    \begin{center}
      \begin{tikzcd}
        0 \ar[r] &
        (\mathrm{F}^s\pi_{k,k+s}(\nu X))[\tau^r] \ar[r] \ar[d] &
        \mathrm{F}^s\pi_{k,k+s}(\nu X) \ar[r, "\cdot \tau^r"] \ar[d, two heads] &
        \mathrm{F}^s\pi_{k,k+s-r}(\nu X) \ar[r] \ar[d, two heads] &
        0 \\
        0 \ar[r] &
        \mathrm{E}_\infty^{s,k,s}[\tau^r] \ar[r] \ar[d, equal] &
        \mathrm{E}_\infty^{s,k,s} \ar[r] \ar[d, equal] &
        \mathrm{E}_\infty^{s,k,s-r} \ar[r] \ar[d, equal] &
        0 \\
        0 \ar[r] &
        \mathrm{B}_{r+1}^{s,k} \ar[r] &
        \mathrm{Z}_\infty^{s,k} \ar[r] &
        \mathrm{Z}_\infty^{s,k}/\mathrm{B}_{r+1}^{s,k} \ar[r] &
        0,
      \end{tikzcd}
    \end{center}
    where the rows are exact by Corollaries \ref{cor:synth-ASS} and \ref{cor:tau-surj}.
    It will suffice to show that left most vertical map is surjective.
    This follows from the snake lemma together with the fact that the map
    $$ \mathrm{F}^{s+1}\pi_{k,k+s}(\nu X) \xrightarrow{\cdot \tau^r} \mathrm{F}^{s+1}\pi_{k,k+s-r}(\nu X) $$
    is surjective, which follows from Corollary \ref{cor:tau-surj}.
    
    (3b) We now suppose that we are given $\alpha \in \pi_k (X)$ detected by $x$. In particular, $x$ is not the target of a differential in the $E$-based Adams spectral sequence for $X$. Then we may modify $\wt{x}$ by elements of higher $\nu E$-filtration without affecting conditions (1a) through (2b). Let $\beta = \tau^{-1} \wt{x}$. Then $\alpha - \beta \in \mathrm{F}^{s+1} \pi_k (X)$. It follows from \Cref{lemm:adams-fil1} that there exists some $e_1 \in \pi_{k,k+s+1} (\nu X)$ such that $\tau^{-1} e_1 = \alpha-\beta$. It follows from \Cref{cor:synth-ASS} that $e_1$ must be in $\nu E$-Adams filtration at least $s+1$. Replacing $\wt{x}$ with $\wt{x} + e_1$, we obtain $\tau^{-1} \wt{x} = \alpha$, as desired.

    (4) Finally, we verify the generation statement.
    Let $A$ denote the $\Z[\tau]$-submodule of $\pi_{k,*}(\nu X)$ generated by the $\wt{x}$, and let $B$ denote the $\tau$-adic completion of $A$. Our first claim is that $B$ remains a natural submodule of $\pi_{k,*}(\nu X)$, which follows from the fact that the $\tau$-adic filtration on $\pi_{k,*}(\nu X)$ is complete and Hausdorff by strong convergence.
		Now, since the inclusion $B \to \pi_{k,*}(\nu X)$ is one between $\tau$-complete objects, we need only note that the map
            $$ B/\tau \to \pi_{k,k+*}(\nu X)/\tau \cong \mathrm{F}^*\pi_{k,k+*}(\nu X)/\mathrm{F}^{*+1}\pi_{k,k+*}(\nu X) \cong \mathrm{E}_\infty^{*,k,*} $$
    is a surjection. The middle isomorphism above follows from \Cref{cor:tau-surj}.
		%
		%From the identification of $\mathrm{F}^\bullet$ with the $\tau$-adic filtration from \Cref{cor:tau-surj} we know that $B$ is $\tau$-adically complete and Hausdorff. Thus, in order to conclude it will suffice to note that by construction the map 
    %$$ B/\tau \to \pi_{k,k+s}(\nu X)/\tau \cong \mathrm{F}^s\pi_{k,k+s}(\nu X)/\tau \cong \mathrm{E}_\infty^{s,k,s} $$
    %is an isomorphism.
\end{proof}

%%% Local Variables:
%%% mode: latex
%%% TeX-master: "Main"
%%% eval: (visual-line-mode t)
%%% eval: (auto-fill-mode 0)
%%% End:
